% CLAIM: at equal population the cheap perturbation aligns with the descent direction
% about as well, and on Countdown it ends at the same accuracy, with full rank ahead early.
% Evidence: tb6_e16.py (E15, E16), phase0/plot.py (E1), plot_e13.py and plot_e19.py
% (E13, E19).
\section{Is the cheaper update as good?}
\label{sec:quality}

A faster update is only worth having if it points in a useful direction. We compare the
perturbations at the same population, on the same loss, two ways: the alignment of one
update with the gradient-descent direction of a differentiable loss, and training on a
task scored by a non-differentiable reward.

\subsection{Alignment with the descent direction}

% F7 -- experiments/countdown/figures/f7-e13-heldout.png (plot_e13.py)
\begin{figure*}[t]
  \centering
  \includegraphics[width=0.84\textwidth]{f7-e13-heldout.png}
  \caption{Countdown with Qwen2.5-0.5B: the share of 2000 held-out puzzles solved,
    against scored training completions (left; 240 per update) and against cumulative
    update time (right), mean over three seeds with a min--max band. The time axis is
    each run's own seconds per update, as each arm was configured, without the first two
    updates (which compile) and without held-out decoding. At 12k samples full rank
    solves 4.4\% against 1.5--2.2\% for ranks 1, 4 and 16, with no overlap between
    seeds; every arm ends near 6\%.}
  \label{fig:f7}
\end{figure*}

We measure alignment as $\cos(\Delta\theta, -\nabla L)$, the cosine between one ES
update and the negative gradient of the next-token loss of Section~\ref{sec:setup}, on
Qwen2.5-0.5B with five perturbation seeds per configuration, and repeat two
configurations on a second prompt batch and on Qwen2.5-1.5B (Table~\ref{tab:tb6}). At
equal population the rank barely matters: rank 1 aligns 0.92--0.98 times as well as full
rank at both model sizes and on both batches, rank 16 0.97--1.05, and rank 4 1.04--1.31
(the 1.31 is at $N{=}30$, where five seeds spread widely: full rank measures 1.3--1.9
and rank 4 0.7--2.0, Table~\ref{tab:tb6-audit}). Unpaired rank 1 is 1.36$\times$
mirrored rank 1 at $N{=}240$, close to the $\sqrt2$ that $N$ against $N/2$ directions
gives, and 1.17$\times$ at $N{=}30$; mirroring halves the directions and in exchange
halves the stored factors.

% generated by experiments/countdown/tb6_e16.py --latex; do not edit
\begin{table}[t]
\centering
\small
\setlength{\tabcolsep}{3.2pt}
\caption{Cosine between one ES update and the exact gradient of the
  next-token loss ($\times 10^{-4}$, median over 5 perturbation seeds), at equal
  population; \emph{1 unp.} is rank 1 without mirrored pairs.
  \emph{ref.}: the mirrored full-rank fit from the synthetic block
  (Figure~\ref{fig:f5}) at the same $N/P$, a post hoc reference; unpaired
  sampling should sit about $\sqrt2$ above it. Batch 2 repeats two
  configurations on freshly drawn prompts. Ranges and the preregistered
  predictions are in Table~\ref{tab:tb6-audit}.}
\label{tab:tb6}
\begin{tabular}{lrrrrrrr}
\toprule
 & & \multicolumn{5}{c}{rank} & \\
\cmidrule(lr){3-7}
model & $N$ & full & 1 & 4 & 16 & 1 unp. & ref. \\
\midrule
0.5B & 30 & $1.42$ & $1.37$ & $1.87$ & $1.49$ & $1.60$ & $1.56$ \\
0.5B & 240 & $4.23$ & $3.91$ & $4.42$ & $4.09$ & $5.33$ & $4.48$ \\
0.5B, batch 2 & 240 & $4.20$ & $3.91$ & -- & -- & -- & $4.48$ \\
1.5B & 30 & $0.83$ & $0.77$ & -- & -- & -- & $0.87$ \\
1.5B & 240 & $2.32$ & $2.27$ & -- & -- & -- & $2.51$ \\
\bottomrule
\end{tabular}
\end{table}


The cosines are of order $10^{-4}$, and that is expected. A Gaussian-projection
approximation predicts alignment proportional to $\sqrt{N/P}$: $\sqrt{3/\pi}\,\sqrt{N/P}$
for unpaired sampling and $1/\sqrt2$ of that for mirrored pairs, and with $P$ = 494M and
$N$ = 30, $\sqrt{N/P}$ is $2.5\times10^{-4}$. The approximation closely describes the
synthetic block, where every arm is within 0.91--1.02 of it at every rank
(Figure~\ref{fig:f5}), and the mirrored arms here measure 0.80--0.92 of it (1.10 for
rank 4 at $N{=}30$). Its applicability to low-rank perturbations depends on the
gradient's structure (Appendix~\ref{sec:alignment}). Most of each update is not along the
descent direction.

The predictions we committed before these runs overestimated low-rank alignment: the
measurements came out at 0.43--0.67 of the prediction at 0.5B, and 23--26\% below it at
1.5B even after a correction fixed in advance. They read each arm's synthetic fit against a
different variable. The $N/P$ reading used here, and Table~\ref{tab:tb6}'s reference
column, were developed afterward; Appendix~\ref{sec:alignment} gives both predictions and
why the first missed.

\subsection{Training on a task}

We trained on Countdown, arithmetic puzzles in the style of TinyZero \citep{tinyzero2025},
with settings adapted from \citet{qiu2025}: their population of 30, noise scale
$\sigma{=}10^{-3}$ and learning rate, converted to Eq.~\ref{eq:update}; centered ranks
where they standardize rewards (the weights' standard deviation is 0.29 against 1, and
we did not match the resulting update magnitudes); and 8 puzzles per update where their
protocol uses 200 samples
(Appendix~\ref{sec:countdown} has the full setup). Each arm runs 500 updates, 120,000
scored training completions, with three seeds, on one A100. A completion scores 1 if it
solves the puzzle and 0.1 if it is well formed but wrong. Held-out accuracy is measured
on 2000 disjoint puzzles by greedy decoding every 50 updates.

The base model solves 1.2\% of the held-out puzzles and formats 43\% of its answers.
Formatting is learned at once: every arm formats 98.5\% or more by the first evaluation,
so Figure~\ref{fig:f7} shows the solve rate, which carries everything that differs
between arms. Training raises it to a plateau, 6.0--6.2\% for full rank and ranks 1, 4
and 16 at the end, and 5.55\% with the embedding frozen (27\% of the parameters). The
approach differs. At the first evaluation, 12k samples in, full rank solves 4.4\% of the
puzzles and ranks 1, 4 and 16 solve 1.5--2.2\%, with the seed ranges apart; by 36k
samples every arm is at 5.3--5.9\%. With one evaluation on the rising part of the curve
we can say that full rank learns faster per sample early on, not by how much.

Per update, rank 1 is faster. As configured, a steady-state update took 4.44\,s at full
rank and 2.62\,s at rank 1, 1.69$\times$. Part of that is how the two arms evaluate: the
full-rank arm scored its members in chunks of five, and the low-rank arm scores the whole
population in one batch. With both batching every member, on one host, the times are
3.62\,s and 2.52\,s: rank 1 updates 1.44$\times$ faster for the same evaluation. Against
update time (Figure~\ref{fig:f7}, right) the early gap mostly closes: full rank solves
4.4\% after 212\,s of updates and 5.5\% after 434\,s, rank 1 5.3\% after 257\,s. The
evaluations are too sparse to say which reaches a given solve rate first, and that axis
uses the arms as configured, not with matched batching.

The task separates training from failure, not small differences in update quality. A
population-16 run, fixed in advance as a positive control because by the square-root scaling
it aligns 27\% worse per update, does not test that either: on about the same budget it
takes 940 updates (120,320 completions) to population 30's 500 and reaches the plateau
sooner (Figure~\ref{fig:f7c}). Among the runs that update all parameters, no rank breaks
training or ends at a different accuracy, on one task at one model size with three
seeds; the run with the embedding frozen ends lower.
